% =============================================================================
% CompetitiveExam/CS502/08-cpu-scheduling-basics-questions.tex
% Competitive Exam Practice 8 — GATE (Computer Science) Pattern
% CPU Scheduling: Concepts, Criteria, FCFS and SJF
%
% Student-facing — NO answers. Answer key: 08-cpu-scheduling-basics-answers.tex.
% Provenance: ALL questions are Original (GATE-pattern). No item in this set is
% attributed to a past GATE paper, and no year/question number is claimed.
% Every numerical answer was computed with a verified scheduling simulator
% (FCFS, non-preemptive SJF, SRTF); no ties occur in any instance.
%
% Source of truth: Slides/CS502/08-cpu-scheduling-basics.tex and
% Notes/CS502/08-cpu-scheduling-basics-notes.tex.
%
% Compile with: cd CompetitiveExam/CS502 &&
%   TEXINPUTS=../../Preambles:$TEXINPUTS xelatex -interaction=nonstopmode 08-cpu-scheduling-basics-questions.tex
% =============================================================================

\documentclass[11pt]{article}
\usepackage[margin=1in]{geometry}
\input{../../Preambles/header}
\coursecode{CS502}
\renewcommand{\thesection}{8.\arabic{section}}

\title{Competitive Exam Practice 8: CPU Scheduling\\\large GATE (Computer Science) Pattern}
\author{Auqib Hamid Lone\\[0.3em]
  \emph{Department of Computer Science \& Engineering}, \emph{GCET Ganderbal}}
\date{\today}

\begin{document}
\maketitle

\noindent Each question carries 1--2 marks. MCQ = single correct option; NAT = Numerical Answer Type. Unless stated otherwise, ignore context-switch overhead, ties are broken in favour of the earlier arrival, and SRTF preempts only when a newcomer's burst is strictly shorter than the running process's remaining time. \emph{All questions in this set are original, written in the GATE pattern; none is claimed to be a past-year question.}

\begin{enumerate}
\item[\textbf{Q1}] \textit{[Original, GATE-pattern]} \textbf{[MCQ, 1 mark]} Consider the transitions (I) running $\to$ waiting, (II) running $\to$ ready, (III) waiting $\to$ ready, (IV) running $\to$ terminated. A \emph{non-preemptive} CPU scheduler makes scheduling decisions only on:
\begin{enumerate}
  \item[(A)] I and IV only \quad (B) II and III only \quad (C) I, II and IV only \quad (D) I, II, III and IV
\end{enumerate}

\item[\textbf{Q2}] \textit{[Original, GATE-pattern]} \textbf{[MCQ, 1 mark]} For a set of processes that are all available at time 0, which policy is guaranteed to give the minimum average waiting time?
\begin{enumerate}
  \item[(A)] FCFS \quad (B) SJF \quad (C) Longest-Job-First \quad (D) Any policy, since total work is fixed
\end{enumerate}

\item[\textbf{Q3}] \textit{[Original, GATE-pattern]} \textbf{[MCQ, 1 mark]} Many I/O-bound processes repeatedly queue behind one CPU-bound process, leaving the I/O devices and then the CPU idle in turn. This behaviour, and the policy that suffers from it, are:
\begin{enumerate}
  \item[(A)] starvation; SJF \quad (B) convoy effect; FCFS \quad (C) thrashing; FCFS \quad (D) convoy effect; SRTF
\end{enumerate}

\item[\textbf{Q4}] \textit{[Original, GATE-pattern]} \textbf{[MCQ, 2 marks]} The next CPU burst is predicted by $\hat{\tau}_{n+1} = \alpha\tau_n + (1-\alpha)\hat{\tau}_n$. Which statement is \textbf{TRUE}?
\begin{enumerate}
  \item[(A)] With $\alpha = 0$, the prediction equals the most recent measured burst
  \item[(B)] With $\alpha = 1$, the prediction never changes from $\hat{\tau}_0$
  \item[(C)] With $0 < \alpha < 1$, the weight on the burst $j$ steps back is $\alpha(1-\alpha)^j$
  \item[(D)] The predictor makes SJF starvation-free
\end{enumerate}

\item[\textbf{Q5}] \textit{[Original, GATE-pattern]} \textbf{[NAT, 2 marks]} Four processes are scheduled FCFS:
\begin{center}\small\begin{tabular}{lcc}\toprule Process & Arrival time (ms) & CPU burst (ms) \\ \midrule \texttt{P1} & 0 & 5 \\ \texttt{P2} & 1 & 3 \\ \texttt{P3} & 3 & 6 \\ \texttt{P4} & 5 & 1 \\ \bottomrule\end{tabular}\end{center}
The average waiting time (in ms) is \rule{1.5cm}{0.4pt}.

\item[\textbf{Q6}] \textit{[Original, GATE-pattern]} \textbf{[NAT, 2 marks]} Four processes are scheduled by non-preemptive SJF:
\begin{center}\small\begin{tabular}{lcc}\toprule Process & Arrival time (ms) & CPU burst (ms) \\ \midrule \texttt{P1} & 0 & 8 \\ \texttt{P2} & 2 & 5 \\ \texttt{P3} & 3 & 2 \\ \texttt{P4} & 4 & 3 \\ \bottomrule\end{tabular}\end{center}
The average turnaround time (in ms) is \rule{1.5cm}{0.4pt}.

\item[\textbf{Q7}] \textit{[Original, GATE-pattern]} \textbf{[NAT, 2 marks]} Four processes are scheduled by SRTF (preemptive SJF):
\begin{center}\small\begin{tabular}{lcc}\toprule Process & Arrival time (ms) & CPU burst (ms) \\ \midrule \texttt{P1} & 0 & 10 \\ \texttt{P2} & 2 & 3 \\ \texttt{P3} & 3 & 1 \\ \texttt{P4} & 6 & 4 \\ \bottomrule\end{tabular}\end{center}
The average waiting time (in ms) is \rule{1.5cm}{0.4pt}.

\item[\textbf{Q8}] \textit{[Original, GATE-pattern]} \textbf{[NAT, 1 mark]} For the SRTF schedule of Q7, the number of times the CPU is switched from one process to a \emph{different} process (not counting the initial dispatch at time 0) is \rule{1.5cm}{0.4pt}.

\item[\textbf{Q9}] \textit{[Original, GATE-pattern]} \textbf{[NAT, 1 mark]} A scheduler uses exponential averaging with $\alpha = 0.5$ and initial guess $\hat{\tau}_0 = 12$~ms. The first two measured CPU bursts are 8~ms and 6~ms. The predicted length (in ms) of the third burst, $\hat{\tau}_2$, is \rule{1.5cm}{0.4pt}.

\item[\textbf{Q10}] \textit{[Original, GATE-pattern]} \textbf{[NAT, 2 marks]} Three processes with CPU bursts 10, 2 and 4~ms all arrive at time 0 and are run non-preemptively in some order. The difference between the largest and the smallest possible average waiting time (in ms, rounded to two decimal places) is \rule{1.5cm}{0.4pt}.
\end{enumerate}

\end{document}
